\documentclass[11pt]{article}

\usepackage[utf8]{inputenc}
\usepackage{amsmath, amssymb}
\usepackage{graphicx}
\usepackage{geometry}
\geometry{margin=1in}

\title{Research on Retinal Disease Classification Using Deep Learning}
\author{Your Name}
\date{\today}

\begin{document}

\maketitle

\begin{abstract}
Retinal diseases are among the leading causes of vision impairment worldwide. Early detection plays a crucial role in preventing irreversible damage. In this paper, we explore the use of deep learning techniques for automated retinal disease classification using image-based datasets.
\end{abstract}

\section{Introduction}
Medical image analysis has significantly benefited from advances in deep learning. In particular, convolutional neural networks (CNNs) have shown strong performance in classification tasks involving retinal images. This study aims to develop an efficient and accurate model for detecting retinal abnormalities.

\section{Related Work}
Previous studies have applied various architectures such as ResNet, VGG, and EfficientNet for medical image classification. Many of these approaches demonstrate high accuracy but often require large computational resources.

\section{Methodology}
The proposed system consists of preprocessing, feature extraction, and classification stages. Images are resized and normalized before being fed into a convolutional neural network model. Data augmentation is applied to improve generalization.

\section{Dataset}
The dataset used in this study contains retinal images labeled into multiple disease categories. Each image is annotated by medical experts to ensure label accuracy.

\section{Results and Discussion}
The model is evaluated using accuracy, precision, recall, and F1-score. Preliminary results show promising performance in distinguishing between healthy and diseased retinal images.

\section{Conclusion}
This study demonstrates the potential of deep learning in assisting early diagnosis of retinal diseases. Future work will focus on improving model robustness and expanding dataset diversity.

\section*{Acknowledgments}
We thank all contributors and dataset providers who made this research possible.

\end{document}